%%%%%%%%%%%%%%%%%%%%%%%%%%%%%%%%%%%%%%%%%%%%%%%%%%%%%%%%%%%%%%%%%%%%%
%
% LateX Template of Björn Schröder
% v0.1
%
%%%%%%%%%%%%%%%%%%%%%%%%%%%%%%%%%%%%%%%%%%%%%%%%%%%%%%%%%%%%%%%%%%%%%


%%%%%%%%%%%%%%%%%%%%%%%%%%%%%%%%%%%%%%%%%%%%%%%%%%%%%%%%%%%%%%%%%%%%%
%%%%% Standard Praeamble %%%%%%%%%%%%%%%%%%%%%%%%%%%%%%%%%%%%%%%%%%%%
%%%%%%%%%%%%%%%%%%%%%%%%%%%%%%%%%%%%%%%%%%%%%%%%%%%%%%%%%%%%%%%%%%%%%

%\documentclass{tufte-book}
\documentclass[12pt,a4paper,openright
    %twocolumn
    ]{memoir}

\usepackage[T1]{fontenc}
\usepackage[utf8]{inputenc}
\usepackage[ngerman]{babel} % Hyphenation patterns if we wanted German
    \renewcommand{\listtablename}{bla} 
    \renewcommand{\listfigurename}{bla} 
    \renewcommand{\contentsname}{bla}

%%%%%%%%%%%%%%%%%%%%%%%%%%%%%%%%%%%%%%%%%%%%%%%%%%%%%%%%%%%%%%%%%%%%%
%%%% My Standard packages
%%%%%%%%%%%%%%%%%%%%%%%%%%%%%%%%%%%%%%%%%%%%%%%%%%%%%%%%%%%%%%%%%%%%%
\usepackage{booktabs}
\usepackage{graphicx}
\usepackage{tabularx}
\usepackage{microtype} 
\usepackage[dvipsnames]{xcolor}
\renewcommand\colorchapnum{\color{ared}}
\renewcommand\colorchaptitle{\color{ared}}
    \definecolor{ared}{rgb}{.647,.129,.149}
    \definecolor{titlepagecolor}{cmyk}{1,.60,0,.40}
    \definecolor{namecolor}{cmyk}{1,.50,0,.10} 
    \definecolor{mygreen}{rgb}{0,0.6,0}
    \definecolor{mygray}{gray}{0.95}
    \definecolor{mymauve}{rgb}{0.58,0,0.82}
\usepackage{colortbl}
%%%%%%%%%%%%%%%%%%%%%%%%%%%%%%%%%%%%%%%%%%%%%%%%%%%%%%%%%%%%%%%%%%%%%
%%%% Useful links and references
%%%%%%%%%%%%%%%%%%%%%%%%%%%%%%%%%%%%%%%%%%%%%%%%%%%%%%%%%%%%%%%%%%%%%
% http://tex.stackexchange.com/questions/19264/techniques-and-packages-to-keep-up-with-good-practices
% http://detexify.kirelabs.org/classify.html
% http://tex.stackexchange.com/questions/553/what-packages-do-people-load-by-default-in-latex

%%%%%%%%%%%%%%%%%%%%%%%%%%%%%%%%%%%%%%%%%%%%%%%%%%%%%%%%%%%%%%%%%%%%%
%%%%%% List of useful packages and associated usecases
%%%%%% Comment if you dont need them
%%%%%%%%%%%%%%%%%%%%%%%%%%%%%%%%%%%%%%%%%%%%%%%%%%%%%%%%%%%%%%%%%%%%%


\RequirePackage[l2tabu,orthodox]{nag}% Old habits die hard. All the same, there are
                                     % commands, classes and packages which are 
                                     % outdated and superseded. nag provides routines
                                     % to warn the user about the use of those.
%\usepackage{amssymb}
%\usepackage{yfonts}
%\usepackage{sidecap}
%\usepackage{textcomp}
%\usepackage{mathptmx}
%\usepackage{courier}
%\usepackage{latexsym}
%\usepackage{amsmath}
%\usepackage{amstext}

\usepackage{listings}
\lstset{ %
    backgroundcolor=\color{mygray},  % choose the background color
    basicstyle=\footnotesize,        % the size of the fonts that are used for the code
    breakatwhitespace=false,         % sets if automatic breaks should only happen at 
                                     % whitespace
    breaklines=true,                 % sets automatic line breaking
    captionpos=b,                    % sets the caption-position to bottom
    commentstyle=\color{mygreen},    % comment style
    deletekeywords={...},            % if you want to delete keywords from the given language
    escapeinside={\%*}{*)},          % if you want to add LaTeX within your code
    extendedchars=true,              % lets you use non-ASCII characters; for 8-bits 
                                     % encodings only, does not work with UTF-8
    frame=single,                    % adds a frame around the code
    keepspaces=true,                 % keeps spaces in text, useful for keeping indentation 
                                     % of code (possibly needs columns=flexible)
    language=Python,                 % the language of the code
    morekeywords={*,...},            % if you want to add more keywords to the set
    numbers=left,                    % where to put the line-numbers; possible values are 
                                     % (none, left, right)
    numbersep=5pt,                   % how far the line-numbers are from the code
    numberstyle=\tiny\color{mygray}, % the style that is used for the line-numbers
    rulecolor=\color{black},         % if not set, the frame-color may be changed on 
                                     % line-breaks within not-black text (e.g. comments 
                                     % (green here))
    showspaces=false,                % show spaces everywhere adding particular underscores; 
                                     % it overrides 'showstringspaces'
    showstringspaces=false,          % underline spaces within strings only
    showtabs=false,                  % show tabs within strings adding particular underscores
    stepnumber=2,                    % the step between two line-numbers. If it's 1, each 
                                     % line will be numbered
    stringstyle=\color{mymauve},     % string literal style
    tabsize=2,                       % sets default tabsize to 2 spaces
    title=\lstname                   % show the filename of files included with 
                                     % \lstinputlisting; also try caption instead of title
}

\newcommand{\incl}[1]{\colorbox{mygray}{\lstinline{#1}} } 
\newcommand{\code}[1]{\colorbox{mygray}{{\texttt{#1}}}}

\usepackage[style=apa,backend=biber]{biblatex} % APA citations
%% I used to use this:
%\usepackage[sort&compress]{natbib}
%    \bibpunct{[}{]}{,}{s}{}{}
%    \bibliographystyle{Bjoern}

\usepackage[some]{background}
\usepackage{geometry}
\usepackage{xcolor}
\usepackage{blindtext}
\usepackage{pifont}
\usepackage{acronym}
\usepackage[textstyle,squaren]{SIunits}
%    \addunit{\M}{\textsc{m}}
%    \addunit{\CV}{\textsc{cv}}
%    \addunit{\dalton}{\textsc{D}a}
%    \addunit{\mol}{mol}
\usepackage{abstract}
\usepackage{ifthen}
%\usepackage{picins}
\usepackage{url}
\usepackage{fancybox}
\usepackage{array}
\usepackage{longtable}
\usepackage{floatflt}
\usepackage{rotating}
\usepackage{lastpage}
\usepackage{currvita}
\usepackage{python}
\usepackage[small,nooneline,bf]{caption}

\usepackage{todonotes}
    %\todo{Rewrite this answer \ldots}
    %\listoftodos

\usepackage{hyperref}
    \hypersetup{
        pdftoolbar=true,        % show Acrobat’s toolbar?
        pdfmenubar=true,        % show Acrobat’s menu?
        pdffitwindow=false,     % window fit to page when opened
        pdfstartview={FitH},    % fits the width of the page to the window
        pdftitle={title},       % title
        pdfauthor={Björn Schröder},     % author
        pdfsubject={Subject},   % subject of the document
        pdfcreator={Björn Schröder},   % creator of the document
        pdfproducer={Björn Schröder}, % producer of the document
        pdfnewwindow=true,      % links in new window
        colorlinks=true,       % false: boxed links; true: colored links
        % color of internal links (change box color with linkbordercolor)
        linkcolor=ared,          
        citecolor=ared,        % color of links to bibliography
        filecolor=ared,      % color of file links
        urlcolor=black           % color of external links
        }


%%%%%%%%%%%%%%%%%%%%%%%%%%%%%%%%%%%%%%%%%%%%%%%%%%%%%%%%%%%%%%%%%%%%%
% Coverpage
%%%%%%%%%%%%%%%%%%%%%%%%%%%%%%%%%%%%%%%%%%%%%%%%%%%%%%%%%%%%%%%%%%%%%

\backgroundsetup{
    scale=1,
    angle=0,
    opacity=1,
    contents={\begin{tikzpicture}[remember picture,overlay]
     \path [fill=titlepagecolor] (current page.west)rectangle (current page.north east); 
     \draw [color=white, very thick] (5,0)--(5,0.5\paperheight);
    \end{tikzpicture}}
    }

\makeatletter                   
\def\printauthor{%                  
    {\large \@author}}          
\makeatother

\author{%
    The Author
    }


%%%%%%%%%%%%%%%%%%%%%%%%%%%%%%%%%%%%%%%%%%%%%%%%%%%%%%%%%%%%%%%%%%%%%
%  Some chapterstyle alternatives for memoir class
%%%%%%%%%%%%%%%%%%%%%%%%%%%%%%%%%%%%%%%%%%%%%%%%%%%%%%%%%%%%%%%%%%%%%

\chapterstyle{hangnum}
\chapterstyle{companion}
\chapterstyle{demo}
\chapterstyle{brotherton}
\chapterstyle{bianchi}
\chapterstyle{dash}
\chapterstyle{ell}
\chapterstyle{demo3}
\chapterstyle{pedersen}


%%%%%%%%%%%%%%%%%%%%%%%%%%%%%%%%%%%%%%%%%%%%%%%%%%%%%%%%%%%%%%%%%%%%%
%  Further customizing
%%%%%%%%%%%%%%%%%%%%%%%%%%%%%%%%%%%%%%%%%%%%%%%%%%%%%%%%%%%%%%%%%%%%%
\widowpenalty100000
\brokenpenalty100000

\definecolor{lightgray}{gray}{0.90}
\newcommand\cb[1]{\fcolorbox{lightgray}{lightgray}{#1}}

% nested items symbol
\renewcommand{\labelitemi}{\tiny{\ding{110}}} 


%%%%%%%%%%%%%%%%%%%%%%%%%%%%%%%%%%%%%%%%%%%%%%%%%%%%%%%%%%%%%%%%%%%%%
%  Modify the margin notes to flushleft
%%%%%%%%%%%%%%%%%%%%%%%%%%%%%%%%%%%%%%%%%%%%%%%%%%%%%%%%%%%%%%%%%%%%%
\marginparmargin{outer}
\let\oldmarginpar\marginpar
\renewcommand\marginpar[2][]{%
    \oldmarginpar{\mpjustification #2}}
%%%%%%%%%%%%%%%%%%%%%%%%%%%%%%%%%%%%%%%%%%%%%%%%%%%%%%%%%%%%%%%%%%%%%

\renewcommand{\rmdefault}{lsbj}

\begin{document}


%%%%%%%%%%%%%%%%%%%%%%%%%%%%%%%%%%%%%%%%%%%%%%%%%%%%%%%%%%%%%%%%%%%%%
% Titlepages 
%%%%%%%%%%%%%%%%%%%%%%%%%%%%%%%%%%%%%%%%%%%%%%%%%%%%%%%%%%%%%%%%%%%%%
%\frontmatter
\begin{titlingpage}
\newgeometry{left=6.5cm} %defines the geometry for the titlepage
\pagecolor{titlepagecolor} % defined in main document

\noindent
\includegraphics[width=6cm]{dragonlogo}\\[-1em]

\color{white}
\noindent
\makebox[0pt][l]{\rule{1.5\textwidth}{1pt}}
\par

\vspace{5mm}
\noindent
{\HUGE{\textsf{Author Name}} \textcolor{namecolor}{\textsf{Institute}}}
\vfill

\noindent
{\huge \textsf{Manual 0.1}}
\vskip\baselineskip

\noindent
\textsf{\today}
\end{titlingpage}
\restoregeometry % restores the geometry
\nopagecolor % Use this to restore the color pages to white


% Die eidesstattliche Erklärung mit Unterschrift
\chapter*{Examplechapter}
\blindtext
\vspace{4cm}
\hspace{2cm} Where, When \hfill Signiture \hspace{2cm}

\newpage
\vfill
\begin{quotation}
\emph{\blindtext}
\end{quotation}
\begin{flushright}
    Dr.\ Who
\end{flushright}
\newpage

\chapter{Another Chapter}
\blindtext
\newpage


\newpage
\tableofcontents*
\newpage
\listoffigures
\newpage
\listoftables
\newpage
\mainmatter

%%%%%%%%%%%%%%%%%%%%%%%%%%%%%%%%%%%%%%%%%%%%%%%%%%%%%%%%%%%%%%%%%%%%%
% Put child docuemtes here
%%%%%%%%%%%%%%%%%%%%%%%%%%%%%%%%%%%%%%%%%%%%%%%%%%%%%%%%%%%%%%%%%%%%%
\blindtext
\blindmathtrue
%%%%%%%%%%%%%%%%%%%%%%%%%%%%%%%%%%%%%%%%%%%%%%%%%%%%%%%%%%%%%%%%%%%%
%%%% Microtypographc Rulez 
%%%%%%%%%%%%%%%%%%%%%%%l%%%%%%%%%%%%%%%%%%%%%%%%%%%%%%%%%%%%%%%%%%%%
\begin{itemize}
\item Kursive Schrift wird generell für Variablen benutzt, aber niemals für Zahlen oder Symbole.
\item \verb+---+ ist ein Gedankenstrich ---
\item \code{$-2$} Minuszeichen werden in Matheumgebung gesetzt
\item \verb+--+ ist ein 'von ... bis' Zeichen
\item \verb+\ldots+ sind \ldots
\item \verb+\,+ ist ein kleiner Abstand in z.\,B.
\item \verb+~+ ist ein geschütztes Lehrzeichen: 
\item \verb+.\+ ist ein Punkt in einer Abkürzung
\end{itemize}



%%%%%%%%%%%%%%%%%%%%%%%%%%%%%%%%%%%%%%%%%%%%%%%%%%%%%%%%%%%%%%%%%%%%
%%  Sammlung nützlicher Befehle  %%%%%%%%%%%%%%%%%%%%%%%%%%%%%%%%%%%
%%%%%%%%%%%%%%%%%%%%%%%%%%%%%%%%%%%%%%%%%%%%%%%%%%%%%%%%%%%%%%%%%%%%

%%%%%%%%%%%%%%%%%%%%%%%%%%%%%%%%%%%%%%%%%%%%%%%%%%%%%%%%%%%%%%%%%%%%%
%%%%%% Listings
%%%%%%%%%%%%%%%%%%%%%%%%%%%%%%%%%%%%%%%%%%%%%%%%%%%%%%%%%%%%%%%%%%%%%
Example for nested lists:
\begin{itemize}
    \item Patent Literature
    \item Scientific Literature
    \item Competitor Information
        \begin{itemize} 
        \item Patent Strategy
        \item New Product Launches
        \end{itemize}
    \item Customer Information
\end{itemize}


%%%%%%%%%%%%%%%%%%%%%%%%%%%%%%%%%%%%%%%%%%%%%%%%%%%%%%%%%%%%%%%%%%%%
% Randnotizen 
%%%%%%%%%%%%%%%%%%%%%%%%%%%%%%%%%%%%%%%%%%%%%%%%%%%%%%%%%%%%%%%%%%%%
\newcommand{\note}[1]{\marginpar[]{\emph{\large{#1}}}}
\note{Randnotiz}


%%%%%%%%%%%%%%%%%%%%%%%%%%%%%%%%%%%%%%%%%%%%%%%%%%%%%%%%%%%%%%%%%%%%
% Bild am Anfang des Absatzes vom Text umflossen 
%%%%%%%%%%%%%%%%%%%%%%%%%%%%%%%%%%%%%%%%%%%%%%%%%%%%%%%%%%%%%%%%%%%%

%\parpic[r]{\includegraphics[width=0.3\textwidth]{dragonlogo}} 
%\blindtext


%%%%%%%%%%%%%%%%%%%%%%%%%%%%%%%%%%%%%%%%%%%%%%%%%%%%%%%%%%%%%%%%%%%%
% Bild in normaler Umgebung
%%%%%%%%%%%%%%%%%%%%%%%%%%%%%%%%%%%%%%%%%%%%%%%%%%%%%%%%%%%%%%%%%%%%
\begin{figure}[!htbp]
\begin{center}
\includegraphics[width=\textwidth]{dragonlogo.pdf}
\end{center}
\caption{\label{}}    
\end{figure}


%%%%%%%%%%%%%%%%%%%%%%%%%%%%%%%%%%%%%%%%%%%%%%%%%%%%%%%%%%%%%%%%%%%%
% Einheiten mit Packet SIUnits 
%%%%%%%%%%%%%%%%%%%%%%%%%%%%%%%%%%%%%%%%%%%%%%%%%%%%%%%%%%%%%%%%%%%%
\begin{itemize}
    \item \unit{170}{\cubic\metre}
    \item power{10}{6} 
    \item \unit{Zellen}{\milli\reciprocal\litre} 
    \item \unit{1000}{$\cdot g$} 
    \item \unit{$-20$}{\celsius}
    \item \unit{100}{\micro\gram\milli\reciprocal\litre}
    \item \unit{109,42}{\degree}
\end{itemize}



%%%%%%%%%%%%%%%%%%%%%%%%%%%%%%%%%%%%%%%%%%%%%%%%%%%%%%%%%%%%%%%%%%%%
% Spezielle Befehle für Sonderzeichen
%%%%%%%%%%%%%%%%%%%%%%%%%%%%%%%%%%%%%%%%%%%%%%%%%%%%%%%%%%%%%%%%%%%%
\begin{itemize}
    \item $\times$
    \item $\mbox{M\textsubscript{r}}=400000$
    \item \emph{aqua polymount}\textregistered
    \item Whatman\texttrademark{}
    \item \textsc{Dynabeads}\textsuperscript{\textregistered{}}
    \item $\sim$\unit{1}{\milli\gram\per\milli\litre}
    \item \textsc{Schrödinger}-Gleichung
\end{itemize}


%%%%%%%%%%%%%%%%%%%%%%%%%%%%%%%%%%%%%%%%%%%%%%%%%%%%%%%%%%%%%%%%%%%%
%  Beispiele für verschiedene Mathe-Umgebungen
%%%%%%%%%%%%%%%%%%%%%%%%%%%%%%%%%%%%%%%%%%%%%%%%%%%%%%%%%%%%%%%%%%%%
\begin{eqnarray*}
E & = & mc^2	\\
E & = &	1/2mv^2
\end{eqnarray*}

\begin{equation}
E=mc^2
\end{equation}

\begin{displaymath}
U_{total}=\sum_{3N+1}~(E_{b}+E_{\theta}+E_{\phi}
\end{displaymath}



%%%%%%%%%%%%%%%%%%%%%%%%%%%%%%%%%%%%%%%%%%%%%%%%%%%%%%%%%%%%%%%%%%%%
% Aufzählungen
%%%%%%%%%%%%%%%%%%%%%%%%%%%%%%%%%%%%%%%%%%%%%%%%%%%%%%%%%%%%%%%%%%%%
\begin{enumerate}
\item Startbedingungen festlegen.
\item Kräfte berechnen.
\end{enumerate}

\begin{itemize}
\item Entwicklung von Mutanten
\item Vorhersage antigener Epitope
\end{itemize}


\blindtext
%%%%%%%%%%%%%%%%%%%%%%%%%%%%%%%%%%%%%%%%%%%%%%%%%%%%%%%%%%%%%%%%%%%%
%  Tabellen  %%%%%%%%%%%%%%%%%%%%%%%%%%%%%%%%%%%%%%%%%%%%%%%%%%%%%%%
%%%%%%%%%%%%%%%%%%%%%%%%%%%%%%%%%%%%%%%%%%%%%%%%%%%%%%%%%%%%%%%%%%%%


%%%%%%%%%%%%%%%%%%%%%%%%%%%%%%%%%%%%%%%%%%%%%%%%%%%%%%%%%%%%%%%%%%%%
% Custom column definitions
%%%%%%%%%%%%%%%%%%%%%%%%%%%%%%%%%%%%%%%%%%%%%%%%%%%%%%%%%%%%%%%%%%%%
\definecolor{Gray}{gray}{0.9}
\newcolumntype{t}{>{\columncolor{Gray}[\tabcolsep]}l}
\newcolumntype{v}[1]{%
>{\begin{turn}{90}
\begin{minipage}{#1}
\raggedright\hspace{-3pt}}l%
<{\end{minipage}
\end{turn}}}

\newcolumntype{g}{>{\columncolor{Gray}[\tabcolsep]}c}
\newcolumntype{b}{>{\columncolor{CornflowerBlue}[\tabcolsep]}c}

%%%%%%%%%%%%%%%%%%%%%%%%%%%%%%%%%%%%%%%%%%%%%%%%%%%%%%%%%%%%%%%%%%%%
% Mehrere Tabellen in einer Umgebung
%%%%%%%%%%%%%%%%%%%%%%%%%%%%%%%%%%%%%%%%%%%%%%%%%%%%%%%%%%%%%%%%%%%%
\begin{longtable}[c]{ll}
\end{longtable}
\begin{table}
\begin{center}
\begin{tabularx}{.8\textwidth}{cllXX}
\toprule 
    This & is & a & long & table \\
\midrule
    This & is & a & long & table \\
    This & is & a & long & table \\
\bottomrule
\end{tabularx}
\end{center}
\caption{\label{}}
\end{table}


%%%%%%%%%%%%%%%%%%%%%%%%%%%%%%%%%%%%%%%%%%%%%%%%%%%%%%%%%%%%%%%%%%%%
% Gedrehte Tabelle mit variabler, angepasster Spaltenbreite siehe Tabularx-Packet
%%%%%%%%%%%%%%%%%%%%%%%%%%%%%%%%%%%%%%%%%%%%%%%%%%%%%%%%%%%%%%%%%%%%
\begin{sidewaystable}
\centering
\begin{tabularx}{\textwidth}{
    >{\hsize=.8\hsize}Xp{50mm}>{\hsize=.8\hsize}X>{\hsize=1.2\hsize}XX>{\hsize=1.2\hsize}X}
\toprule  
    This & is & a & very & nice &table  \\
\midrule
    This & is & a & very & nice &table  \\
    This & is & a & very & nice &table  \\
\bottomrule
\end{tabularx} 
\caption{\label{}}
\end{sidewaystable}






\blindtext
\end{document}
